% !Mode:: "TeX:UTF-8"

\chapter{基于度量学习的视觉 3D 人体姿态估计方法(测试)}

\section{引言(测试)}
本研究针对现有人体姿态估计模型在实际应用中因环境复杂多变导致的泛化性不足问题，尝试了基于度量学习的视觉 3D 人体姿态估计方法。传统方法通过增加模型复杂度或扩充数据集提升指标性能，但面临两个核心矛盾：(1) 3D 标注数据获取成本高昂；(2)预训练模型难以适应新环境条件。本研究的创新性在于，基于已训练完成的 2D 姿态估计模型，通过引入新环境参考图像构建域适应机制，使模型无需重新训练即可快速适应不同光照、背景等环境变量，直接输出适配当前场景的 2D 姿态结果。最终结合三角测量方法实现 3D 关节坐标推算，本章的整体网络架构如图 3-1 所示。本章在实验环节，通过消融实验证明了参考视图个数对 2D 估计结果的影响，以及引入的图卷积结构的有效性，并在两个数据集中对比了本论文方法与其他方法的性能。

\section{基于度量学习的多视角 2D 人体姿态估计(测试)}
本节将介绍 2D 姿态估计环节的各个模块，首先介绍用于初步提取关节特征的网络，为了在 2D 阶段充分利用图像中的语义信息，本文选用 Swin Transformer[79]作为特征提取网络，分别提取待估计视图和参考视图中的关节特征。其次介绍相似性度量网络及其注意力热图的表现形式，最后介绍图卷积结构和特征解码结构。本节的 2D 姿态估计网络结构如图所示，在本节的最后，会给出该网络架构对应的实验结果与分析。

\subsection{特征提取网络(测试)}
随着 Alexnet 卷积神经网络在图像领域取得重大突破后，各类卷积神经网络也呈现爆发式增长，在人体姿态估计领域中，由于人体目标的尺度不一样，为了让模型能分辨不同尺度的人体目标，有些研究人员也考虑到对不同分辨率的人体图像做处理，提升模型的鲁棒性，使其能适应不同尺度的人体目标图像。

\subsection{特征编码结构(测试)}
随着 Alexnet 卷积神经网络在图像领域取得重大突破后，各类卷积神经网络也呈现爆发式增长，在人体姿态估计领域中，由于人体目标的尺度不一样，为了让模型能分辨不同尺度的人体目标，有些研究人员也考虑到对不同分辨率的人体图像做处理，提升模型的鲁棒性，使其能适应不同尺度的人体目标图像。

\section{基于三角测量方法的 3D 人体关节坐标求解(测试)}
本节将介绍 2D 姿态估计环节的各个模块，首先介绍用于初步提取关节特征的网络，为了在 2D 阶段充分利用图像中的语义信息，本文选用 Swin Transformer[79]作为特征提取网络，分别提取待估计视图和参考视图中的关节特征。其次介绍相似性度量网络及其注意力热图的表现形式，最后介绍图卷积结构和特征解码结构。本节的 2D 姿态估计网络结构如图所示，在本节的最后，会给出该网络架构对应的实验结果与分析。

\subsection{对极几何理论(测试)}
随着 Alexnet 卷积神经网络在图像领域取得重大突破后，各类卷积神经网络也呈现爆发式增长，在人体姿态估计领域中，由于人体目标的尺度不一样，为了让模型能分辨不同尺度的人体目标，有些研究人员也考虑到对不同分辨率的人体图像做处理，提升模型的鲁棒性，使其能适应不同尺度的人体目标图像。

\subsection{3D关节坐标求解(测试)}
随着 Alexnet 卷积神经网络在图像领域取得重大突破后，各类卷积神经网络也呈现爆发式增长，在人体姿态估计领域中，由于人体目标的尺度不一样，为了让模型能分辨不同尺度的人体目标，有些研究人员也考虑到对不同分辨率的人体图像做处理，提升模型的鲁棒性，使其能适应不同尺度的人体目标图像。